\subsection{局部凸空间的完备化}

定理 1（Grothendieck）. --- 设 E 是一个局部凸空间，并设 S 是 E 上一个适配且覆盖的有界结构。设 \(F \subset E^{*}\) 是 E 上这样的线性型所成的空间：它们在每个属于 S 的集上的限制是连续的。若给 F 赋予 S-拓扑，则从 \(E'_{S}\) 到 F 中的典范单射扩张为从完备化 \(\hat{E}'_{S}\)（即 \(E'_{S}\) 的完备化）到 F 上的一个同构。

由于 E 上线性型的每个简单极限是线性型（III, 第 16 页, 命题 4），且 E 上的有界结构 \(\mathfrak{S}\) 是覆盖的，由 GT, X, § 1, No.~6, 推论 2 可得：带 \(\mathfrak{S}\) -拓扑的空间 F 是 Hausdorff 且完备的。显然 \(E_{\mathfrak{S}}'\) 是 F 的一个拓扑向量子空间；因此只需证明 \(E_{\mathfrak{S}}'\) 在 F 中处处稠密。这由下面的引理得出：

引理 1. --- 设 A 是 E 的一个闭凸均衡子集，并设 u 是 E 上的一个线性型，它在 A 上的限制连续。则对每个 \(\varepsilon > 0\)，存在一个 \(x' \in E'\) 使得

\[
\left| u (x) - \langle x, x ^ {\prime} \rangle \right| \leqslant \varepsilon \quad \text { for   every } \quad x \in A.
\]

设 \(\varepsilon > 0\)。存在 E 中 0 的一个闭凸均衡邻域 U，使得 \(|u(x)| \leqslant \varepsilon\) 对所有 \(x \in \mathrm{U} \cap \mathrm{A}\) 成立。我们知道极 \(\mathrm{U}^{\circ}\)（U 在 \(\mathrm{E}^{*}\) 中的）含于 \(\mathrm{E}'\) 且对拓扑 \(\sigma(\mathrm{E}^{*}, \mathrm{E})\) 是紧的（III, 第 17 页, 推论 2）。由于极 \(\mathrm{A}^{\circ}\)（A 在 \(\mathrm{E}^{*}\) 中的）对 \(\sigma(\mathrm{E}^{*}, \mathrm{E})\) 是闭的，由此得出 \(\mathrm{A}^{\circ} + \mathrm{U}^{\circ}\) 是 \(\mathrm{E}^{*}\) 的一个闭凸子集（GT, III, § 4, No.~1, 推论 1）。

设 C 是 E 的一个闭凸均衡子集。则 C 对 \(\sigma(\mathrm{E},\mathrm{E}^{\prime})\) 是闭的（II, 第 45 页, 推论 3），从而对 \(\sigma(\mathrm{E},\mathrm{E}^{*})\) 也是闭的，因而我们有 \(C = C^{\circ\circ}\)（对 E 与 \(E^{*}\) 之间的对偶而言）。结果，我们有

\[
\mathrm{A} \cap \mathrm{U} = \mathrm{A} ^ {\circ \circ} \cap \mathrm{U} ^ {\circ \circ} = (\mathrm{A} ^ {\circ} \cup \mathrm{U} ^ {\circ}) ^ {\circ} \supset (\mathrm{A} ^ {\circ} + \mathrm{U} ^ {\circ}) ^ {\circ}
\]

由此得到

\[
(\mathrm{A} \cap \mathrm{U}) ^ {\circ} \subset (\mathrm{A} ^ {\circ} + \mathrm{U} ^ {\circ}) ^ {\circ \circ} = \mathrm{A} ^ {\circ} + \mathrm{U} ^ {\circ}.
\]

由于线性型 \(\varepsilon^{-1}u\) 属于 \((\mathrm{A}\cap\mathrm{U})^{\circ}\)，存在 \(v\in A^{\circ}\) 及 \(w\in U^{\circ}\) 使得 \(u=\varepsilon(v+w)\)。于是 \(x'=\varepsilon w\) 属于 \(E'\)，且 \(u-x'=\varepsilon v\) 在 A 上的绝对值以 \(\varepsilon\) 为上界；引理得证。

现在设 E 是一个局部凸 Hausdorff 空间，\(\hat{E}\) 是它的完备化。E 上每个连续线性型 f 由连续性扩张到 \(\hat{E}\)：因而我们有 \((\hat{\mathrm{E}})' = \mathrm{E}'\)（III, 第 16 页），且 \(\hat{\mathbf{E}}\) 的每个元素定义了 \(\mathbf{E}'\) 上的一个线性型；也就是说，定义了代数对偶 \(\mathbf{E}^{\prime *}\)（\(\mathbf{E}'\) 的）的一个元素。此外，\(\mathbf{E}\)（相应地，\(\hat{\mathbf{E}}\)）与 \(\mathbf{E}'\) 之间的对偶是分离的（II, 第 24 页, 推论 1）。因此 \(\mathbf{E}\) 与 \(\hat{\mathbf{E}}\) 可与 \(\mathbf{E}^{\prime *}\) 的向量子空间等同起来。

定理 2. --- 设 E 是一个局部凸 Hausdorff 空间，\(\hat{\mathbf{E}}\) 是它的完备化；我们把 E 与 \(\hat{\mathbf{E}}\) 与 \(\mathbf{E}^{\prime *}\) 的向量子空间等同起来。则为使元素 \(f \in \mathbf{E}^{\prime *}\) 属于 \(\hat{\mathbf{E}}\)，必要且充分的是：\(f\) 在 \(\mathbf{E}'\) 的每个等度连续子集上的限制对拓扑 \(\sigma(\mathbf{E}', \mathbf{E})\) 连续。

空间 E 可与 \(E'\) 的拓扑对偶等同起来，当给 \(E'\) 赋予拓扑 \(\sigma(E', E)\) 时（II, 第 43 页, 命题 3）；另一方面，若 S 是 \(E'\) 的等度连续子集所成之集，则 E 上给定的拓扑就是 S-拓扑（III, 第 19 页, 推论 1）。于是由 III, 第 13 页, 命题 1 可得，S 的诸集对 \(\sigma(E', E)\) 有界（参见后面 III, 第 22 页, 命题 9）；换言之，S 是对拓扑 \(\sigma(E', E)\) 而言适配且覆盖的有界结构。若把 E 换成 \(E'\) 并把 \(E'_{S}\) 换成 E，则定理 2 是定理 1 的推论。

推论 1（Banach）. --- 设 E 是一个 Hausdorff 且完备的局部凸空间。为使 \(E'\) 上的一个线性型对弱拓扑 \(\sigma(E', E)\) 连续（即由 E 的一个元素产生），只需它在 \(E'\) 的每个等度连续子集上的限制对 \(\sigma(E', E)\) 连续。

评注. --- 此外假设 E 中存在一个可数完全集；则 \(E'\) 的每个等度连续子集对拓扑 \(\sigma(E', E)\) 可度量化（III, 第 18 页, 命题 6）；因此为验证 \(E'\) 上的线性型 u 是弱连续的，只需验证：对每个等度连续序列 \((x_{n}')\)（在 \(E'\) 中且对 \(\sigma(E', E)\) 收敛于 0），我们有 \(\lim_{n \to \infty} u(x_{n}') = 0\)。

推论 2. --- 设 \((\mathrm{E}_{i})_{i\in\mathbb{I}}\) 是 Hausdorff 局部凸空间的一个族，并设 E 是它们的拓扑直和。则从诸 \(\hat{E}_{i}\) 的直和到 \(\hat{E}\) 中的典范映射是一个同构。特别地，E 完备当且仅当所有 \(E_{i}\) 都完备。

我们知道 E 的对偶可与诸 \(E_{i}\) 的对偶的乘积等同起来（II, 第 30 页, 公式 (1)）。设 \(u \in \hat{E}\)，并设 \(u_{i} \in E_{i}^{\prime *}\) 是 u（视为 \(E^{\prime *}\) 的元素）在 \(E_{i}^{\prime} \subset E^{\prime}\) 上的限制。立即可见，只需证明：\(u_{i} = 0\) 除有限多个指标 \(i \in I\) 外成立。反之，假设存在一个由互异指标构成的序列 \((i_{n})_{n \in \mathbf{N}}\) 使得 \(u_{i_{n}} \neq 0\)。则存在 \(x_{i_{n}} \in E_{i_{n}}^{\prime}\) 使得 \(u_{i_{n}}(x_{i_{n}}) = n\)。所有 \(x_{i_{n}}\) 所成之集 H 在 \(E^{\prime}\) 中等度连续，而 u 在 H 上的限制不是有界的，这是不可能的。

